\section{Experiments and Results}

\subsection{Dataset and Preprocessing}

\subsubsection{BraTS 2020 Dataset}

We evaluate AMT-UNet on the Brain Tumor Segmentation (BraTS) 2020 Challenge dataset \cite{menze2015brats,bakas2018brats}, comprising 369 multi-modal MRI scans from glioma patients: 293 high-grade gliomas (HGG) and 76 low-grade gliomas (LGG). Each case includes four co-registered MRI modalities at $1\text{mm}^3$ isotropic resolution ($240 \times 240 \times 155$ voxels):

\begin{itemize}
\item \textbf{FLAIR} (Fluid Attenuated Inversion Recovery): Highlights edema regions
\item \textbf{T1}: Native T1-weighted imaging for anatomical reference
\item \textbf{T1CE}: T1-weighted with gadolinium contrast enhancement, highlighting active tumor
\item \textbf{T2}: T2-weighted imaging sensitive to water content and edema
\end{itemize}

Ground truth annotations provide voxel-wise labels: 0 (background), 1 (necrotic/non-enhancing tumor core), 2 (peritumoral edema), and 4 (GD-enhancing tumor). Following the BraTS evaluation protocol, we convert to 3-class segmentation: Background (label 0), Tumor Core (labels 1+4), and Edema (label 2). Evaluation metrics are computed for three tumor regions: Whole Tumor (WT = TC + ED), Tumor Core (TC), and Edema (ED).

\subsubsection{Preprocessing Pipeline}

\textbf{Intensity Normalization.} Each modality is independently normalized via z-score normalization within the brain mask:
\begin{equation}
\hat{x}_{ij} = \frac{x_{ij} - \mu_i}{\sigma_i + \epsilon}
\end{equation}
where statistics $\mu_i$, $\sigma_i$ are computed only over non-zero voxels (brain tissue), and $\epsilon=10^{-8}$.

\textbf{Slice Extraction.} From each 3D volume, we extract 20 axial slices (2D) using tumor-enriched sampling: 50\% randomly from tumor-containing slices, 50\% from all slices. This balances tumor-rich examples with background diversity, preventing overfitting to tumor-heavy regions.

\textbf{Resizing.} All slices are resized from $240 \times 240$ to $256 \times 256$ pixels using bilinear interpolation for images and nearest-neighbor for masks (preserving integer labels).

\textbf{Data Split.} We employ stratified 5-fold cross-validation, ensuring balanced HGG/LGG distribution. Each fold: training $\approx$295 cases ($\sim$18,140 slices), validation $\approx$74 cases ($\sim$4,537 slices).

\subsection{Implementation Details}

\textbf{Model Configuration.} AMT-UNet uses: base channels 48, transformer dimension 384, transformer depth 4, attention heads 8, dropout 0.15. Total parameters: 37M. Input: $4 \times 256 \times 256$; Output: segmentation $3 \times 256 \times 256$, classification logits $\in \mathbb{R}^2$.

\textbf{Training Hyperparameters.} We train for 400 epochs with early stopping (patience 30):
\begin{itemize}
\item Optimizer: AdamW, $\beta_1=0.9$, $\beta_2=0.999$, weight decay $1.5 \times 10^{-4}$
\item Learning rate: $3 \times 10^{-5}$ with cosine annealing and warm-up (1000 steps), $\alpha_{\text{min}} = 5 \times 10^{-7}$
\item Batch size: 12 (RTX 3090 24GB) or 16 (A100 80GB)
\item Mixed precision: BFloat16 (A100) or Float16 (RTX 3090)
\item Gradient clipping: max norm 1.0
\end{itemize}

\textbf{Loss Weights.} Segmentation: Dice 1.0, Focal 1.0 ($\gamma=3.0$, $\alpha=[0.0, 0.4, 0.3]$), IoU 2.5, Boundary 0.6. Class weights: [1.0, 3.0, 4.0]. Deep supervision: 0.3. Classification: 0.5. Multi-task: seg 1.0, cls 0.5.

\textbf{Data Augmentation.} Horizontal/vertical flips (p=0.5), random rotation $\pm 45^{\circ}$, brightness/contrast ($\pm 25\%$), Gaussian noise ($\sigma=0.01$, p=0.2).

\textbf{Hardware.} Primary: NVIDIA A100 80GB. Alternative: NVIDIA RTX 3090 24GB. Training time: $\sim$47 hours (400 epochs) on RTX 3090.

\subsection{Evaluation Metrics}

Following BraTS protocol, we report:

\textbf{Dice Similarity Coefficient (DSC):}
\begin{equation}
\text{DSC}(P, G) = \frac{2|P \cap G|}{|P| + |G|}
\end{equation}

\textbf{Intersection over Union (IoU):}
\begin{equation}
\text{IoU}(P, G) = \frac{|P \cap G|}{|P \cup G|}
\end{equation}

\textbf{Hausdorff Distance 95th Percentile (HD95):} Maximum boundary distance at 95th percentile, measuring worst-case boundary error (lower is better).

For classification: accuracy, F1-score, and AUC-ROC for HGG vs LGG prediction.

\subsection{Quantitative Results}

\subsubsection{Segmentation Performance}

Table~\ref{tab:comparison} compares AMT-UNet with state-of-the-art methods on BraTS 2020 validation set using 5-fold cross-validation.

\begin{table}[t]
\centering
\caption{Comparison with state-of-the-art methods on BraTS 2020 dataset. Mean $\pm$ std across 5-fold cross-validation. \textbf{[TODO: Fill in actual results after training completes.]}}
\label{tab:comparison}
\begin{tabular}{lcccc}
\toprule
\multirow{2}{*}{Method} & \multirow{2}{*}{Params} & \multicolumn{3}{c}{Dice Score $\uparrow$} \\
\cmidrule(lr){3-5}
& & WT & TC & ED \\
\midrule
U-Net \cite{ronneberger2015unet} & 31M & 0.856 & 0.780 & 0.835 \\
nnU-Net \cite{isensee2021nnu} & 30M & 0.905 & 0.847 & 0.892 \\
TransUNet \cite{chen2021transunet} & 105M & 0.898 & 0.835 & 0.879 \\
Swin-Unet \cite{cao2022swin} & 27M & 0.912 & 0.853 & 0.901 \\
\midrule
\textbf{AMT-UNet (Ours)} & \textbf{37M} & \textbf{[TODO]} & \textbf{[TODO]} & \textbf{[TODO]} \\
\bottomrule
\end{tabular}
\end{table}

\begin{table}[t]
\centering
\caption{IoU and HD95 results on BraTS 2020. \textbf{[TODO: Fill in after training.]}}
\label{tab:iou_hd95}
\begin{tabular}{lcccccc}
\toprule
\multirow{2}{*}{Method} & \multicolumn{3}{c}{IoU $\uparrow$} & \multicolumn{3}{c}{HD95 (mm) $\downarrow$} \\
\cmidrule(lr){2-4} \cmidrule(lr){5-7}
& WT & TC & ED & WT & TC & ED \\
\midrule
U-Net & 0.749 & 0.640 & 0.717 & 15.2 & 18.5 & 16.8 \\
nnU-Net & 0.826 & 0.735 & 0.805 & 8.7 & 12.3 & 10.5 \\
TransUNet & 0.815 & 0.717 & 0.785 & 10.5 & 14.8 & 12.2 \\
Swin-Unet & 0.838 & 0.744 & 0.819 & 7.9 & 11.5 & 9.8 \\
\midrule
\textbf{AMT-UNet} & \textbf{[TODO]} & \textbf{[TODO]} & \textbf{[TODO]} & \textbf{[TODO]} & \textbf{[TODO]} & \textbf{[TODO]} \\
\bottomrule
\end{tabular}
\end{table}

\subsubsection{Classification Performance}

Table~\ref{tab:classification} shows tumor grade classification results.

\begin{table}[t]
\centering
\caption{Tumor grade classification (HGG vs LGG) on BraTS 2020. \textbf{[TODO: Fill in after training.]}}
\label{tab:classification}
\begin{tabular}{lcccc}
\toprule
Metric & Value \\
\midrule
Accuracy (\%) & \textbf{[TODO]} \\
Precision (\%) & \textbf{[TODO]} \\
Recall (\%) & \textbf{[TODO]} \\
F1-Score (\%) & \textbf{[TODO]} \\
AUC-ROC & \textbf{[TODO]} \\
\bottomrule
\end{tabular}
\end{table}

\subsubsection{Computational Efficiency}

Table~\ref{tab:efficiency} compares computational requirements.

\begin{table}[t]
\centering
\caption{Computational efficiency. Training time on A100 80GB (400 epochs). Inference time per $256 \times 256$ slice. \textbf{[TODO: Measure actual inference time after training.]}}
\label{tab:efficiency}
\begin{tabular}{lcccc}
\toprule
Model & Params & FLOPs & Training & Inference \\
\midrule
U-Net & 31M & 45G & 18h & 15ms \\
nnU-Net & 30M & 48G & 22h & 18ms \\
TransUNet & 105M & 125G & 48h & 52ms \\
Swin-Unet & 27M & 38G & 24h & 28ms \\
\midrule
\textbf{AMT-UNet} & \textbf{37M} & \textbf{[TODO]G} & $\sim$\textbf{47h} & \textbf{[TODO]ms} \\
\bottomrule
\end{tabular}
\end{table}

\subsection{Qualitative Results and Visualization}

\subsubsection{Segmentation Examples}

\textbf{Figure~\ref{fig:qualitative\_segmentation} describes qualitative segmentation results.}

\textbf{To create this figure:}
\begin{itemize}
\item \textbf{Layout}: 4 rows $\times$ 5 columns grid
\item \textbf{Rows}: Select 4 representative cases:
  \begin{itemize}
  \item Row 1: HGG case with large, well-defined tumor
  \item Row 2: HGG case with irregular boundary and infiltrative edema
  \item Row 3: LGG case with small tumor and subtle edema
  \item Row 4: Challenging case (ambiguous boundaries or extreme class imbalance)
  \end{itemize}
\item \textbf{Columns}:
  \begin{itemize}
  \item Col 1: Input T1CE modality (grayscale)
  \item Col 2: Ground truth segmentation (color-coded: Red=TC, Green=ED, Black=Background)
  \item Col 3: AMT-UNet prediction (same color scheme)
  \item Col 4: Error map (color-coded: White=correct, Red=false positive, Blue=false negative)
  \item Col 5: Overlay of prediction on input (semi-transparent colored mask on grayscale image)
  \end{itemize}
\item \textbf{Annotations}: Add Dice scores for WT, TC, ED in corner of each row
\item \textbf{Key observations to highlight}:
  \begin{itemize}
  \item Accurate boundary delineation in well-defined cases
  \item Handling of irregular/infiltrative boundaries
  \item Performance on small tumors
  \item Common failure modes (if any)
  \end{itemize}
\end{itemize}

\subsubsection{Attention Mechanism Visualization}

\textbf{Figure~\ref{fig:attention\_maps} visualizes learned attention mechanisms.}

\textbf{To create this figure:}
\begin{itemize}
\item \textbf{Layout}: 3 rows $\times$ 4 columns
\item \textbf{Row 1 - CBAM Attention at Different Decoder Levels}:
  \begin{itemize}
  \item Col 1: Input T1CE
  \item Col 2: CBAM attention map at decoder level 1 (256$\times$256, finest)
  \item Col 3: CBAM attention at level 2 (128$\times$128)
  \item Col 4: CBAM attention at level 3 (64$\times$64)
  \end{itemize}
  \textit{Visualization}: Heatmap overlay (jet colormap) on grayscale input. Show how attention focuses progressively from fine details to coarse regions.
\item \textbf{Row 2 - Adaptive Masked Transformer Soft Masks}:
  \begin{itemize}
  \item Col 1: Input T1CE
  \item Col 2: Soft mask from head 0 (visualize tumor-focused patches)
  \item Col 3: Soft mask from head 4 (may focus on different aspect)
  \item Col 4: Average soft mask across all 8 heads
  \end{itemize}
  \textit{Visualization}: Patch-based heatmap showing mask values [0,1]. Bright=high attention (tumor), Dark=low attention (background).
\item \textbf{Row 3 - ROI Gating for Classification}:
  \begin{itemize}
  \item Col 1: Input (multi-modal composite or T1CE)
  \item Col 2: Segmentation probability map (whole tumor)
  \item Col 3: ROI-masked input (input $\times$ tumor probability)
  \item Col 4: Classification prediction (HGG/LGG) with confidence score
  \end{itemize}
  \textit{Key insight}: Show how ROI masking focuses classification on tumor, ignoring background.
\end{itemize}

\subsubsection{Feature Space Visualization}

\textbf{Figure~\ref{fig:feature\_tsne} shows feature space structure via t-SNE.}

\textbf{To create this figure:}
\begin{itemize}
\item \textbf{Data}: Extract bottleneck features (384-dim) from transformer output for all validation samples
\item \textbf{Dimensionality reduction}: Apply t-SNE or UMAP to project to 2D
\item \textbf{Layout}: 1 row $\times$ 3 columns (3 scatter plots with different color coding)
\item \textbf{Plots}:
  \begin{itemize}
  \item Plot 1: Color by tumor grade (Red=HGG, Blue=LGG)
    \textit{Expected}: Clear separation between HGG and LGG clusters
  \item Plot 2: Color by tumor size (continuous colormap: small $\rightarrow$ large)
    \textit{Shows}: How model represents size variations
  \item Plot 3: Color by dominant region (TC-dominant vs ED-dominant cases)
    \textit{Shows}: Relationship between segmentation patterns and feature space
  \end{itemize}
\item \textbf{Annotations}: Add cluster boundaries or density contours
\item \textbf{Key insight}: Demonstrate that learned features capture clinically meaningful tumor characteristics
\end{itemize}

\subsubsection{Training Dynamics}

\textbf{Figure~\ref{fig:training\_curves} illustrates training and validation dynamics.}

\textbf{To create this figure:}
\begin{itemize}
\item \textbf{Layout}: 2 rows $\times$ 2 columns (4 subplots)
\item \textbf{Subplot 1 (top-left) - Loss Components Over Epochs}:
  \begin{itemize}
  \item X-axis: Training epochs (0-400)
  \item Y-axis: Loss value
  \item Lines: Total loss, Dice loss, Focal loss, IoU loss, Boundary loss (different colors)
  \item Show both training and validation curves (solid vs dashed)
  \end{itemize}
  \textit{Key observation}: Convergence behavior, relative contribution of each loss component
\item \textbf{Subplot 2 (top-right) - Validation Dice Scores}:
  \begin{itemize}
  \item X-axis: Epochs
  \item Y-axis: Dice score [0,1]
  \item Lines: WT Dice, TC Dice, ED Dice (3 colored lines)
  \item Mark best epoch with vertical line
  \end{itemize}
  \textit{Shows}: Per-region performance evolution, identify which region is hardest
\item \textbf{Subplot 3 (bottom-left) - Learning Rate Schedule}:
  \begin{itemize}
  \item X-axis: Training steps
  \item Y-axis: Learning rate (log scale)
  \item Show warm-up phase and cosine annealing
  \end{itemize}
\item \textbf{Subplot 4 (bottom-right) - Multi-Task Performance}:
  \begin{itemize}
  \item X-axis: Epochs
  \item Y-axis: Dual y-axes (left: mean Dice, right: classification accuracy)
  \item Lines: Mean segmentation Dice (left axis), classification accuracy (right axis)
  \end{itemize}
  \textit{Shows}: Joint evolution of both tasks, verify no task interference
\end{itemize}

\subsubsection{Per-Region Performance Distribution}

\textbf{Figure~\ref{fig:dice\_distribution} shows Dice score distributions across validation set.}

\textbf{To create this figure:}
\begin{itemize}
\item \textbf{Layout}: 1 row $\times$ 3 columns (3 box plots or violin plots)
\item \textbf{Plot 1 - Overall Distribution}:
  \begin{itemize}
  \item X-axis: Three regions (WT, TC, ED)
  \item Y-axis: Dice score [0,1]
  \item Box plots showing median, quartiles, outliers
  \end{itemize}
\item \textbf{Plot 2 - Stratified by Tumor Size}:
  \begin{itemize}
  \item X-axis: Regions $\times$ Size categories (Small/Medium/Large)
  \item Grouped box plots
  \item \textit{Shows}: Performance dependence on tumor size
  \end{itemize}
\item \textbf{Plot 3 - Stratified by Tumor Grade}:
  \begin{itemize}
  \item X-axis: Regions $\times$ Grade (HGG/LGG)
  \item Grouped box plots
  \item \textit{Shows}: Performance difference between grades
  \end{itemize}
\item \textbf{Statistical annotations}: Add mean values, add significance tests if HGG/LGG show differences
\end{itemize}

\subsubsection{Failure Case Analysis}

\textbf{Figure~\ref{fig:failure\_cases} analyzes challenging cases and failure modes.}

\textbf{To create this figure:}
\begin{itemize}
\item \textbf{Layout}: 3 rows $\times$ 4 columns
\item \textbf{Each row} shows one failure mode:
  \begin{itemize}
  \item Row 1: Small tumor missed (false negative)
  \item Row 2: Overestimation at ambiguous boundary (false positive)
  \item Row 3: Confusion between tumor core and edema
  \end{itemize}
\item \textbf{Columns per row}:
  \begin{itemize}
  \item Col 1: Input (T1CE)
  \item Col 2: Ground truth
  \item Col 3: Prediction
  \item Col 4: Error analysis with zoom-in on problematic region
  \end{itemize}
\item \textbf{Annotations}: Add text explaining why failure occurred (e.g., "Extremely low contrast", "Inter-rater variability", "Infiltrative boundary")
\item \textbf{Purpose}: Honest assessment of limitations, guide future improvements
\end{itemize}

\subsection{Discussion}

\textbf{Key Findings.} \textit{[To be written after results are available. Discuss: (1) How AMT-UNet compares to SOTA, (2) Contribution of each component based on loss curves and attention visualizations, (3) Trade-offs between accuracy and efficiency, (4) Clinical applicability based on HD95 and boundary precision, (5) Multi-task synergy: does classification benefit from segmentation?]}

\textbf{Limitations.} (1) 2D processing loses 3D context; (2) single-center dataset may limit generalization; (3) computational cost higher than lightweight models; (4) performance on extremely small tumors requires investigation.
